\documentclass[12pt]{article}
\usepackage[utf8]{utf8}
\usepackage[margin=1.5cm]{geometry}
\usepackage{booktabs}
\usepackage{tabularx}
\usepackage{array}

% Ajusta a espessura da linha interna (fina) do booktabs
\setlength{\lightrulewidth}{0.2pt}

\begin{document}

% --- TÍTULOS CENTRALIZADOS ---
\begin{center}
    {\fontsize{15pt}{18pt}\selectfont \textbf{Teoria do Aprendizado Estatístico / PI 3}}\\[0.2cm]
    {\fontsize{11pt}{14pt}\selectfont \textbf{Professor:} João Paulo Ferreira de Mello}\\[0.1cm]
    {\fontsize{11pt}{14pt}\selectfont \textbf{Integrantes:} Ivan Guimarães, Jorge Sugano, Gustavo}\\[0.3cm]
\end{center}

% --- DESCRIÇÃO DO DATASET ---
\noindent
\fontsize{10pt}{13pt}\selectfont
\textbf{Descrição do Dataset:} O presente estudo utiliza a base de dados do \textbf{Índice Paulista de Desenvolvimento Municipal (IPDM)}, elaborada pela Fundação Seade. Trata-se de um indicador sintético desenvolvido para avaliar e acompanhar o desempenho socioeconômico dos 645 municípios do Estado de São Paulo em periodicidade bianual (2014--2024). A base mensura as condições municipais por meio de um índice geral e de três dimensões fundamentais: Riqueza, Longevidade e Escolaridade, fornecendo insumos para a análise comparativa da evolução do desenvolvimento humano e econômico no estado. Dicionário de dados segue abaixo:

\vspace{0.25cm}

% --- TABELA DO DICIONÁRIO DE DADOS ---
\begin{table}[htbp]
\centering
\fontsize{8pt}{10.5pt}\selectfont % Tamanho ajustado para acomodar o respiro vertical das linhas
\renewcommand{\arraystretch}{1.3} % Dá o espaço vertical para as letras não colarem nas linhas

\caption{Dicionário de Dados --- Índices e Indicadores do IPDM (2014--2024)}
\label{tab:dicionario_ipdm}
\vspace{0.15cm}

\begin{tabularx}{\textwidth}{>{\ttfamily}l l >{\raggedright\arraybackslash}p{4.8cm} X l}
\toprule
\textbf{Variável} & \textbf{Tipo} & \textbf{Domínio / Variação} & \textbf{Descrição} & \textbf{Fonte} \\
\midrule
cod\_ibge    & Numérico & [3500105, 3557303] & Código oficial de 7 dígitos do IBGE para o município & IBGE \\
\midrule
municipio   & Texto    & Nomes de municípios & Nome do município do Estado de São Paulo & IBGE \\
\midrule
valor       & Numérico & [0,000; 1,000] & Nota/pontuação da cidade no IPDM geral ou na dimensão & Seade \\
\midrule
ano         & Numérico & [2014, 2024] & Ano de referência dos dados (série bianual) & Seade \\
\midrule
tipo        & Texto    & Categórico & Escopo da linha (Geral, Riqueza, Longevidade ou Escolaridade) & Seade \\
\midrule
valor\_estado & Numérico & [0,000; 1,000] & Pontuação equivalente para o total do Estado de SP & Seade \\
\midrule
indicador1  & Texto / Num. & ["Consumo anual de energia elétrica residencial (MWh) por ligação", "Taxas de mortalidade infantil (por mil nascidos vivos)", "Taxas de atendimento escolar de crianças de 0 a 3 anos (\%)"] & 1º componente específico da dimensão de análise & Seade \\
\midrule
indicador2  & Texto / Num. & ["Consumo anual de energia elétrica comercial, serviços e rural (MWh) por ligação", "Taxas de mortalidade perinatal (por mil nascidos vivos)", "Proporção média de alunos do 5º ano com proficiência em Língua Portuguesa e Matemática (\%)"] & 2º componente específico da dimensão de análise & Seade \\
\midrule
indicador3  & Texto / Num. & ["Rendimento do trabalho formal mais benefícios previdenciarios per capita (R\$ de 2024)", "Taxas de mortalidade 15 a 39 anos (por mil hab.)", "Proporção média de alunos do 9º ano com proficiência em Língua Portuguesa e Matemática (\%)"] & 3º componente específico da dimensão de análise & Seade \\
\midrule
indicador4  & Texto / Num. & ["Produto Interno Bruto per capita (R\$ de 2024)", "Taxas de mortalidade 60 a 69 anos (por mil hab.)", "Taxas de distorção idade-série no Ensino Médio (\%)"] & 4º componente específico da dimensão de análise & Seade \\
\midrule
indicador5  & Texto / Num. & ["Indicador Riqueza", "Indicador Longevidade", "Indicador Escolaridade"] & Identificador ou nota do indicador sintético da dimensão & Seade \\
\bottomrule
\end{tabularx}
\end{table}

\end{document}